\documentclass[12pt,a4paper]{article}

\usepackage[a4paper,margin=1in]{geometry}
\usepackage{graphicx}
\usepackage{booktabs}
\usepackage{amsmath,amssymb}
\usepackage{float}
\usepackage{hyperref}
\usepackage{xcolor}
\usepackage{listings}
\usepackage{enumitem}
\usepackage{fancyhdr}
\usepackage{longtable}

\hypersetup{colorlinks=true,linkcolor=blue,urlcolor=blue}

\pagestyle{fancy}
\fancyhf{}
\lhead{ICS1512}
\rhead{Machine Learning Algorithms Laboratory}
\cfoot{\thepage}

\lstset{
language=Python,
basicstyle=\ttfamily\small,
numbers=left,
frame=single,
breaklines=true,
showstringspaces=false
}

\begin{document}

\begin{tabular}{|p{4cm}|p{11cm}|}
\hline
Experiment No. & 2 \\ \hline
Experiment Title & K-Nearest Neighbors and Na\"{i}ve Bayes Classification
\footnote{\textbf{GitHub Repository: }\url{https://github.com/Naren-Kandasamy/Machine_Learning_Lab-3122247001036}}\\ \hline
Student Name & Naren Kandasamy \\ \hline
Register Number & 3122247001036 \\ \hline
Faculty & Dr. Poreddy Ajay Kumar Reddy \\ \hline
Submission Date & 02/08/2026 \\ \hline
\end{tabular}

\section{Objective}
The purpose of this experiment was to compare the performance of different variations of the Na\"{i}ve Bayes algorithm (Gaussian, Multinomial, and Bernoulli) against a standard K-Nearest Neighbors (KNN) model. Utilizing the Kaggle Spambase dataset, the experiment also involved tuning the hyperparameter $k$, assessing the computational differences between KDTree and BallTree structures, and employing Cross Validation to ensure the reliability of the results.

\section{Problem Statement}
The primary objective of this task was to classify emails as either "Spam" or "Ham" based on the frequency of specific words and characters appearing in the text. The input dataset consisted of a table containing 57 numerical features (representing word frequencies), while the expected output was a binary prediction where 1 indicates spam and 0 indicates a normal email.

\section{Dataset Description}
\begin{longtable}{|p{6cm}|p{8cm}|}
\hline
Dataset Name & Kaggle Spambase Dataset (\texttt{spambase\_csv.csv}) \\ \hline
Dataset Source & Kaggle / UCI Machine Learning Repository \\ \hline
Number of Samples & 4,601 \\ \hline
Number of Features & 57 (Word/Character Frequencies and lengths) \\ \hline
Number of Classes & 2 (Spam, Ham) \\ \hline
Missing Values & 0 \\ \hline
Train-Test Split & Train: 3,220 | Validation: 690 | Test: 691 \\ \hline
\end{longtable}

\section{Exploratory Data Analysis}
An initial examination of the dataset was conducted to understand its structure. The vocabulary size comprised 56 predictive features alongside a single target label. A class distribution chart was plotted to confirm the balance between spam and ham emails, and it was determined that the dataset was reasonably balanced, allowing the analysis to proceed without extensive resampling.

\section{Data Preprocessing}
Preparing the data for modeling required a few specific processing steps:
\begin{itemize}
\item \textbf{Handling Sparsity:} Given that the data consisted of text frequencies, the standard mean-imputation and zero-mean scaling were omitted for the Na\"{i}ve Bayes models. This decision was made to preserve the exact frequency counts of the words without introducing artificial distortion.
\item \textbf{Scaling for KNN:} Because KNN relies on calculating physical distances between points, it is highly sensitive to unscaled data. A \texttt{MinMaxScaler} was applied to both the training and testing sets prior to fitting the KNN model to ensure fair distance weighting.
\item \textbf{Train-test split:} The dataset was divided into Training, Validation, and Testing sets using a stratified approach, ensuring that the ratio of spam to ham remained consistent across all splits.
\end{itemize}

\section{Machine Learning Algorithm}
Two main types of algorithms were evaluated for this classification task:

\textbf{1. Na\"{i}ve Bayes (Gaussian, Multinomial, Bernoulli):}
\begin{itemize}
\item \textbf{Working principle:} A probabilistic classifier based on Bayes' theorem. It operates on the strict assumption that every feature is completely independent of the others, which, despite being an oversimplification for text data, often yields strong results in practice.
\item \textbf{Mathematical formulation:} $P(c|x) = \frac{P(x|c)P(c)}{P(x)}$ where $c$ represents the class and $x$ represents the feature vector.
\end{itemize}

\textbf{2. K-Nearest Neighbors (KNN):}
\begin{itemize}
\item \textbf{Working principle:} Rather than constructing an underlying probability model, this instance-based algorithm examines the $k$ closest data points in the training set and assigns a class based on a majority vote.
\item \textbf{Mathematical formulation:} Distance is typically calculated using the Euclidean formula: $d(p, q) = \sqrt{\sum_{i=1}^{n} (q_i - p_i)^2}$.
\item \textbf{Tree Structures:} KDTree and BallTree structures were utilized to accelerate the distance computations, as calculating distances to every single point iteratively becomes computationally expensive.
\end{itemize}

\section{Experimental Setup}
\begin{tabular}{ll}
\toprule
Python Version & 3.10+ \\
Scikit-Learn Version & 1.0+ \\
IDE & Jupyter Notebook / VS Code \\
Operating System & Linux \\
Hardware & Standard CPU Workspace \\
\bottomrule
\end{tabular}

\section{Hyperparameter Selection}

\textbf{KNN Comparison (Varying $k$):}
Several different values for $k$ (1, 3, 5, 7, 9, 11) were tested. Prior to the introduction of distance weighting, the highest accuracy observed was at $k=5$ (approximately 90.3\%).

\begin{tabular}{ll}
\toprule
Optimization Strategy & Result \\
\midrule
GridSearchCV Best Params & \texttt{\{'n\_neighbors': 7, 'weights': 'distance'\}} (Found in 0.51s) \\
RandomizedSearchCV Best Params & \texttt{\{'n\_neighbors': 5, 'weights': 'distance'\}} (Found in 0.29s) \\
\bottomrule
\end{tabular}

\section{Results}

\textbf{1. Model Performance Comparison:}
\begin{center}
\begin{tabular}{lccccc}
\toprule
Model & Accuracy & Precision & Recall & F1-score & ROC-AUC \\
\midrule
Multinomial NB & 0.8958 & 0.9464 & 0.7794 & 0.8548 & 0.9558 \\
Gaussian NB & 0.8220 & 0.7018 & 0.9522 & 0.8081 & 0.9166 \\
Bernoulli NB & 0.6483 & 0.8717 & 0.1250 & 0.2186 & 0.6097 \\
Optimized KNN & 0.9088 & 0.9023 & 0.8654 & 0.8835 & 0.9412 \\
\bottomrule
\end{tabular}
\end{center}

\textbf{2. 5-Fold Cross Validation (Best NB vs Best KNN):}
\begin{center}
\begin{tabular}{lcc}
\toprule
Fold & Na\"{i}ve Bayes & Best KNN \\
\midrule
1 & 0.8825 & 0.9055 \\
2 & 0.8964 & 0.9130 \\
3 & 0.8849 & 0.9015 \\
4 & 0.8900 & 0.9105 \\
5 & 0.9015 & 0.9118 \\
\midrule
\textbf{Average} & \textbf{0.8911} & \textbf{0.9085} \\
\bottomrule
\end{tabular}
\end{center}

\textbf{3. Tree Structure Timing Comparison:}
\begin{center}
\begin{tabular}{lcc}
\toprule
Algorithm & Train Time (s) & Predict Time (s) \\
\midrule
KDTree & 0.0107 & 0.1106 \\
BallTree & 0.0071 & 0.0942 \\
\bottomrule
\end{tabular}
\end{center}

\section{Result Visualization}
A series of plots were generated to visualize model performance during execution:
\begin{itemize}
\item \textbf{Accuracy vs k:} A line chart illustrating the fluctuation in accuracy as the value of $k$ varied from 1 to 11.
\item \textbf{ROC Curves \& PR Curves:} These curves facilitated the comparison of the area under the curve across the different Na\"{i}ve Bayes models and the optimized KNN model.
\item \textbf{Confusion Matrices:} Grids were generated to identify where the models were producing false positives versus false negatives.
\item \textbf{GridSearchCV Heatmap:} A heatmap visualization clearly displaying which combinations of $k$ and weight parameters yielded the optimal F1-Score.
\end{itemize}

\section{Discussion}
The results revealed that Multinomial Na\"{i}ve Bayes was highly effective, achieving nearly 90\% test accuracy and an ROC-AUC of 0.9558. However, after tuning the KNN classifier using GridSearchCV (which identified optimal settings of $k=7$ and distance-based weights), the KNN model outperformed Na\"{i}ve Bayes. It achieved a stable average accuracy of 90.85\% across the 5-fold cross-validation. Additionally, a comparison of the two tree algorithms showed that BallTree was marginally faster than KDTree for making predictions on this specific dataset.

\section{Conclusion}
The experiment demonstrated that while probabilistic models such as Na\"{i}ve Bayes are highly computationally efficient and perform well on text data, distance-based algorithms like KNN can yield superior classification performance when the data is appropriately scaled and hyperparameters are carefully tuned. Both models ultimately proved successful at distinguishing between spam and ham emails.

\section{References}
\begin{enumerate}
    \item Pedregosa, F., et al. "Scikit-learn: Machine Learning in Python." \textit{Journal of Machine Learning Research}, 2011.
\end{enumerate}

\end{document}
